% === III. Tujuan dan Manfaat Dikembangkannya Perangkat Lunak ===
\secbox{III.}{Tujuan dan Manfaat Dikembangkannya Perangkat Lunak}

\subsecbar{3.1 Tujuan Pengembangan Perangkat Lunak}
Secara khusus, tujuan pengembangan LUNA meliputi:
\begin{enumerate}[leftmargin=*,label=\arabic*.]
    \item Menyediakan layanan pendampingan psikologis awal yang dapat diakses kapan saja.
    \item Menghasilkan diari emosional secara otomatis melalui \textit{Automatic Diary Generation}.
    \item Mempertahankan konteks percakapan antar-sesi menggunakan \textit{Long-term Conversation Memory}.
    \item Mendeteksi dini risiko krisis emosional melalui \textit{Early Mental Crisis Detection}.
    \item Menyediakan interaksi suara yang natural melalui fitur \textit{Voice Conversation}.
\end{enumerate}

\subsecbar{3.2 Manfaat Pengembangan Perangkat Lunak}

\subsubsecstyle{Bagi Pengguna (Generasi Z)}
\begin{enumerate}[leftmargin=*,label=\arabic*.]
    \item Menyediakan ruang aman untuk mengekspresikan kondisi emosional.
    \item Membantu mengenali pola emosi dan melakukan refleksi diri.
    \item Memberikan pendampingan psikologis yang lebih personal dan berkelanjutan.
\end{enumerate}

\subsubsecstyle{Bagi Masyarakat}
\begin{enumerate}[leftmargin=*,label=\arabic*.]
    \item Mendukung deteksi dini dan pencegahan krisis kesehatan mental.
    \item Meningkatkan literasi kesehatan mental serta mengurangi stigma.
    \item Memperluas akses pendampingan psikologis awal, terutama bagi masyarakat dengan keterbatasan layanan profesional.
\end{enumerate}

\subsubsecstyle{Bagi Pengembangan AI di Bidang Kesehatan}
\begin{enumerate}[leftmargin=*,label=\arabic*.]
    \item Mengimplementasikan arsitektur Multi-Agent AI untuk pendampingan kesehatan mental.
    \item Memanfaatkan Retrieval-Augmented Generation (RAG) guna meningkatkan akurasi dan mengurangi halusinasi informasi.
    \item Mendukung pengembangan AI yang lebih aman melalui mekanisme deteksi krisis dan AI Safety Guardrails.
\end{enumerate}
